% Page 045

\seite{45}

\abschnitt{Ernte}
\pstart
\margmark{In diessiger Gegend\\werden. Jedoch auch\\Drittel und Grünfutt\\frühern Summarisirt\\noch Geraden gebrach\\Witterung zog sich d\\ziemlich wenig auf e\\Bedürfnis fiel im al\\bei einem mittleren}%
kann es allein als noch befriedigend angesehen\ob
die Preise der erste und Roggen aufs\ob
er konnten wegen des Mangels nicht als in\ob
werden, ein grosser Teil wurde davon\ob
t worden. Wegen vielfach ungünst.\ob
ie Grummeternte in die Länge. Obst gab es\ob
ine, was fast die Hälfte verfaulte ganz das\ob
lgemeinen nicht. Die Kartoffeln waren\ob
Ertrag.

Gelesen 14.\ 12.\ 08

\pend
\begin{verse}
1908
\end{verse}
\pstart

Weihnachtsfeier.\ob
\margmark{Am Geburtstag Sr. Ma\\nächst Besuch der Ge\\Seffern mithin Kinde\\geziemte Aufführsten\\Kindern einstudirten\\Kindern vorgetragen\\Kindergeld. Lehrer K}%
j. des Kaisers wurden den\ob
mein eingeladenen Bürgern zu\ob
r des Herrn Ortsschulinspektor Pfarrer\ob
 mittheilen, die von den grössern\ob
, etlichen Gedichte, die von einigen\ob
wurden. Die Ansprachen hielt\ob
espen aus Schleid.

Prüfung.\ob
\margmark{Die diesjähr. Prüfe\\den Herrn Kreisschul}%
wurde am 14.\ Dezember 08 durch\ob
inspektor Herz aus Bitburg erwähnt.

Gelesen 30/3.\ 09
Schmitt, Pfr.

Osterprüfung.\ob
\margmark{Die Osterprüfung wur\\schulinspektor Herrn}%
de am 30.\ März durch den Herrn Orts\ohb
 Pfarrer Schmitt aus Seffern abgehalten.

Das Schuljahr 1909--10.\ob
\margmark{Entlassen wurden in\\den widmen sich der}%
diesem Jahre 5 Knaben. Ausge\ohb
Landwirthschaft. Aufgenommen war
\pend
